\subsection{Homological dimension of a ring}

DEFINITION 2. --- One calls the homological dimension of A and denotes by dh(A) the least upper bound in \(\overline{\mathbf{Z}}\) of the set of integers \(n\) for which there exist two A-modules M and N such that \(\operatorname{Ext}_{\mathbf{A}}^{n}(\mathbf{M},\mathbf{N})\neq0\) .

One has \(\mathrm{dh}(0) = -\infty\) , \(\mathrm{dh(A)} \geqslant 0\) if \(A \neq 0\) . One will see below that \(\mathrm{dh(A)} = 1\) if \(A\) is principal and is not a field, and that, if \(K\) is a commutative field

\[
\mathrm{dh} (\mathrm{K} [ \mathrm{X} _ {1}, \dots , \mathrm{X} _ {n} ]) = n.
\]

PROPOSITION 4. --- Let n be an integer ≥ 0. The following conditions are equivalent: (i) dh(A) ≤ n,

\begin{enumerate}
\def\labelenumi{(\roman{enumi})}
\setcounter{enumi}{1}
\tightlist
\item
  for every A-module M, one has \(\mathrm{dp}_{\mathbf{A}}(\mathbf{M}) \leqslant n\) ,
\end{enumerate}

(ii') for every finitely generated A-module M, one has \(\mathrm{dp}_{\mathbf{A}}(\mathbf{M}) \leqslant n\) ,

\begin{enumerate}
\def\labelenumi{(\roman{enumi})}
\setcounter{enumi}{2}
\tightlist
\item
  for every exact sequence
\end{enumerate}

\[
0 \rightarrow \mathrm{K} \rightarrow \mathrm{P} _ {n - 1} \rightarrow \mathrm{P} _ {n - 2} \rightarrow \dots \rightarrow \mathrm{P} _ {0}
\]

where the \(P_{i}\) are projective, K is projective,

\begin{enumerate}
\def\labelenumi{(\roman{enumi})}
\setcounter{enumi}{3}
\tightlist
\item
  for every exact sequence
\end{enumerate}

\[
\mathrm{I} ^ {0} \rightarrow \mathrm{I} ^ {1} \rightarrow \dots \rightarrow \mathrm{I} ^ {n - 1} \rightarrow \mathrm{N} \rightarrow 0
\]

where the \(\mathrm{I}^{i}\) are injective, N is injective,

\begin{enumerate}
\def\labelenumi{(\alph{enumi})}
\setcounter{enumi}{21}
\tightlist
\item
  every A-module has an injective resolution of length ≤ n.
\end{enumerate}

The equivalence of conditions (i), (ii) and (iii) follows from prop. 1. We obviously have (ii) \(\Rightarrow\) (ii'). It therefore suffices for us to prove (ii') \(\Rightarrow\) (iv) \(\Rightarrow\) (v) \(\Rightarrow\) (i).

(ii') \(\Rightarrow\) (iv): with the notations of (iv), let K be the kernel of \(\mathrm{I}^{0} \to \mathrm{I}^{1}\) . By X, p.~128, cor. 4, we have for every A-module M an isomorphism \(\operatorname{Ext}_{\mathbf{A}}^{1}(\mathbf{M}, \mathbf{N}) \to \operatorname{Ext}_{\mathbf{A}}^{n+1}(\mathbf{M}, \mathbf{K})\) ; it then follows from (ii') that \(\operatorname{Ext}_{\mathbf{A}}^{1}(\mathbf{M}, \mathbf{N})\) is zero for every finitely generated A-module M. By X, p.~93, prop. 11, this implies that N is injective, whence (iv).

\begin{enumerate}
\def\labelenumi{(\roman{enumi})}
\setcounter{enumi}{3}
\tightlist
\item
  \(\Rightarrow\) (v): let M be an A-module. Applying (iv) to the exact sequence
\end{enumerate}

\[
0 \rightarrow \mathrm{M} \rightarrow \mathrm{I} ^ {0} (\mathrm{M}) \rightarrow \mathrm{I} ^ {1} (\mathrm{M}) \rightarrow \dots \rightarrow \mathrm{I} ^ {n - 1} (\mathrm{M}) \rightarrow \mathrm{K} ^ {n - 1} (\mathrm{M}) \rightarrow 0
\]

of X, p.~52, we conclude that \(\mathbf{K}^{n-1}(\mathbf{M})\) is injective, whence (v).

\begin{enumerate}
\def\labelenumi{(\alph{enumi})}
\setcounter{enumi}{21}
\tightlist
\item
  \(\Rightarrow\) (i): this follows from X, p.~100, th. 1.
\end{enumerate}

Remarks. --- 1) If \(\mathrm{dh}(\mathbf{A}) \leqslant n < +\infty\) , we have \(\operatorname{Tor}_{n+1}^{\mathbf{A}}(\mathbf{P}, \mathbf{M}) = 0\) for every A-module M and every right A-module P, since \(\mathrm{dp}_{\mathbf{A}}(\mathbf{M}) \leqslant n\) (cf.~lemma 1).

\begin{enumerate}
\def\labelenumi{\arabic{enumi})}
\setcounter{enumi}{1}
\tightlist
\item
  For dh(A) to be finite, it is necessary and sufficient that dp\_A(M) be finite for every nonzero A-module M. This indeed follows from what precedes and from X, p.~135, cor. 1.
\end{enumerate}

COROLLARY. --- Suppose A is left noetherian and let n be an integer \textgreater{} 0. The following conditions are equivalent:

\begin{enumerate}
\def\labelenumi{(\roman{enumi})}
\item
  \(\mathrm{dh(A)}\leqslant n,\)
\item
  for every pair of finitely generated A-modules M and N, we have \(\operatorname{Ext}_{\mathbf{A}}^{n+1}(\mathbf{M}, \mathbf{N}) = 0\) ,
\item
  for every finitely generated left A-module M and every finitely generated right A-module P, we have \(\operatorname{Tor}_{n+1}^{\mathbf{A}}(\mathbf{P},\mathbf{M}) = 0\) .
\end{enumerate}

This follows from props. 2 and 4.

Remark. --- By the equivalence of (i) and (iii), we have dh(A) = dh(A°) if A is right and left noetherian. This equality is not satisfied in general (X, p.~204, exercise 20).

PROPOSITION 5. --- Suppose A is left noetherian and of finite homological dimension and let \(\mathcal{C}_{0}\) (resp. \(\mathcal{C}\) ) be the set of isomorphism classes of finitely generated projective A-modules (resp. of finitely generated A-modules). Then the canonical homomorphism of Grothendieck groups \(\mathrm{K}(\mathcal{C}_{0}) \to \mathrm{K}(\mathcal{C})\) is bijective.

This follows from X, p.~137, cor.

